% -------------------------------------------------------
\section{Limitations and Future Work}
\label{sec:rs-limitations}

The Recursive Spawning Protocol provides strong correctness guarantees — drift prevention, chain integrity, and termination — under the assumptions of honest initialization, secure cryptographic primitives, and Fact Layer write protocol enforcement. These guarantees come with structural costs that honest engineering requires us to acknowledge. This section examines three specific limitations, each grounded in a named failure mode, and charts directions for future work that follow directly from each.

% -------------------------------------------------------
\subsection{Performance Overhead}
\label{sec:rs-limitations-overhead}

RSP's enforcement mechanisms impose runtime overhead that conventional agentic runtimes do not incur. This overhead arises from three sources: the Fact Layer pre-commit hook, the forensic ledger hash-chaining, and the chain verification procedure at each spawn event.

The Fact Layer pre-commit hook evaluates two conditions — scope subset relationship and depth constraint — for every spawn event. In a conventional runtime, a spawn event requires a single authorization decision by the parent node. Under RSP, the pre-commit hook additionally queries the forensic ledger to confirm the parent's provisioned state, formally verifies the scope subset relationship, and records the spawn event in the ledger before the child node is provisioned. This adds a cryptographic write operation (hash-chaining the new entry) and a scope evaluation to the spawn critical path.

The forensic ledger's hash-chaining provides tamper-evidence but incurs a sequential write dependency: each entry's hash must be computed after the preceding entry's hash is known, precluding parallel write optimization for high-throughput spawn scenarios. In chains with deep nesting and frequent sibling spawn events — typical in high-density task decomposition workloads — this sequential bottleneck may become a throughput constraint.

We characterize the overhead formally. Let $T_{\mathrm{precommit}}$ be the pre-commit hook evaluation time, $T_{\mathrm{ledger\_write}}$ be the ledger append and hash-chain write time, and $T_{\mathrm{verify}}(|C|)$ be the chain verification time at provisioning, which grows with chain depth $|C|$ because verification must traverse the full authorization lineage. The total spawn overhead is:

\begin{equation}
T_{\mathrm{spawn\_overhead}} = T_{\mathrm{precommit}} + T_{\mathrm{ledger\_write}} + T_{\mathrm{verify}}(|C|)
\label{eq:spawn-overhead}
\end{equation}

At moderate chain depths ($k \leq 10$), this overhead is negligible relative to model inference and tool execution times. At deep chains ($k \geq 50$), the verification cost may become non-trivial in latency-sensitive workloads.

Mitigation strategies include ledger checkpointing (periodic pruning of historical entries beyond a fixed depth, retaining hash-chain integrity for recent segments), parallel pre-commit validation for sibling spawn events originating from the same parent, and hardware-accelerated cryptographic primitives for hash computation. These are engineering optimizations, not protocol modifications; they do not alter RSP's correctness properties.

% -------------------------------------------------------
\subsection{Scalability Boundaries}
\label{sec:rs-limitations-scalability}

RSP's correctness guarantees apply to single spawn chains — sequences of nodes connected by the parent pointer relation. The protocol as specified does not address multi-parented agent graphs, lateral agent migration, or dynamic coalition formation among agents with independent human anchors.

In a system where agents from different human principals interact — for example, a multi-tenant platform where each tenant's RSP chain operates on shared execution infrastructure — RSP alone does not govern inter-chain interactions. Each chain's intra-chain integrity is preserved, but cross-chain coordination events are not scoped by any RSP constraint. A compromised agent in one chain could influence the behavior of agents in another chain through shared infrastructure (memory stores, message queues, tool APIs), creating a lateral compromise pathway that RSP's chain-scoped guarantees cannot detect.

The current RSP addresses intra-chain accountability exclusively. Inter-chain isolation is addressed by other \bxthree pillars — Spatial Firewall and Sandbox Gate — which enforce execution compartmentation. RSP's scalability boundary is therefore bounded by the coverage of these complementary pillars. In deployments where inter-chain isolation is critical and Spatial Firewall and Sandbox Gate are not co-deployed, RSP's guarantees do not extend to cross-chain interaction scenarios.

Extending RSP to inter-chain accountability requires an inter-chain interaction protocol that enforces cross-chain scope constraints at coordination events. This extension is co-dependent with Spatial Firewall and Sandbox Gate co-deployment; the combined system would require a unified correctness statement that RSP alone cannot provide.

% -------------------------------------------------------
\subsection{Compromised Purpose Layer}
\label{sec:rs-limitations-purpose-layer}

RSP's correctness properties assume that the Purpose Layer's spawn authorization policy $P_{\mathrm{spawn}}$ is well-formed at provisioning. If the Purpose Layer itself is compromised — if an adversarial actor can modify $P_{\mathrm{spawn}}$ after provisioning, or if the initial purpose clause contains a scope definition that is intentionally over-broad or ambiguously specified — then the drift prevention theorem's base case is undermined.

Specifically, if $R(n_0)$ at the human anchor is specified with insufficient granularity, the scope refinement constraint $R(n_{i+1}) \subseteq R(n_i)$ does not prevent the human anchor's intent from being progressively diluted across generations of agents. An adversarial initial provisioning could set $R(n_0)$ to encompass virtually all operational states, and subsequent refinements — each satisfying the subset constraint — could steer the chain toward a target state that no reasonable interpretation of $R(n_0)$ would endorse. The \emph{structure} of the refinement chain is correct; the \emph{semantic content} of the root authorization is not.

This limitation is not addressable by RSP alone. It is addressed by Purpose Layer governance: rigorous auditing of initial purpose clause definitions, formal review of scope boundaries before provisioning, and ongoing Purpose Layer integrity monitoring. RSP provides the structural enforcement mechanism; it does not validate the semantic correctness of what the Purpose Layer authorizes. The assumption of well-formed provisioning is load-bearing for the correctness properties, and the security of the Purpose Layer provisioning process is a prerequisite — not an optional enhancement — for RSP's guarantees to hold.

% -------------------------------------------------------
\subsection{Future Work}
\label{sec:rs-future-work}

Three directions for future work follow directly from the limitations above.

\medskip
\noindent\textbf{Successor promotion for legitimate chain reconfiguration.} When a parent node is intentionally decommissioned, RSP currently requires escalation to the human anchor for every descendant reassignment — friction that may be unacceptable in high-throughput production systems. A successor promotion mechanism would permit operational continuity by transferring authority from a terminated node to a pre-designated successor, without modifying the immutable parent pointer. The historical chain remains intact in the forensic ledger; only the operational routing is updated via a new ledger entry. This is the highest-priority protocol extension.

\medskip
\noindent\textbf{Parallel ledger writes for high-throughput spawn.} The sequential hash-chain write bottleneck can be addressed by a checkpoint-based two-tier ledger architecture: fine-grained hash chains within checkpoint intervals (enabling parallel write batches) and coarse-grained checkpoint hashes across interval boundaries (maintaining tamper-evidence at the coarse level). The trade-off is increased verification complexity for deep historical queries — acceptable for operational workloads where verification depth is bounded by the checkpoint interval.

\medskip
\noindent\textbf{Purpose Layer integrity monitoring.} Automated provisioning-time validation tools — scope boundary formal verification, ambiguity detection in natural-language purpose specifications, automated test generation for scope refinement chains — would address the compromised Purpose Layer risk directly. This is a deployment-time validation layer that strengthens RSP's assumptions before the protocol's guarantees are operative, not a modification of the protocol itself.
